\documentclass[a4paper,11pt]{article}

\usepackage[T2A]{fontenc}
\usepackage[utf8]{inputenc}
\usepackage[russian]{babel}

\usepackage{amsmath, amssymb}
\usepackage{xcolor}
\usepackage[most]{tcolorbox}
\usepackage{geometry}

\geometry{left=1cm,right=1cm,top=1.2cm,bottom=1.2cm}

\setlength{\parindent}{0pt}
\setlength{\parskip}{3pt}

\setlength{\abovedisplayskip}{4pt}
\setlength{\belowdisplayskip}{4pt}
\setlength{\abovedisplayshortskip}{3pt}
\setlength{\belowdisplayshortskip}{3pt}

\allowdisplaybreaks

% ---------- Цвета ----------
\definecolor{questioncolor}{RGB}{46, 125, 50}
\definecolor{answercolor}{RGB}{25, 118, 210}
\definecolor{definitioncolor}{RGB}{30, 136, 229}
\definecolor{theoremcolor}{RGB}{198, 40, 40}
\definecolor{proofcolor}{RGB}{239, 108, 0}

% ---------- Блоки ----------
\newtcolorbox{questionbox}{
    enhanced,
    breakable,
    colback=questioncolor!5,
    colframe=questioncolor,
    boxrule=0pt,
    leftrule=4pt,
    arc=0mm,
    left=6pt,
    right=6pt,
    top=4pt,
    bottom=4pt,
    before skip=4pt,
    after skip=4pt,
    fontupper=\small
}

\newtcolorbox{answerbox}{
    enhanced,
    breakable,
    colback=answercolor!4,
    colframe=answercolor,
    boxrule=0pt,
    leftrule=4pt,
    arc=0mm,
    left=6pt,
    right=6pt,
    top=4pt,
    bottom=4pt,
    before skip=4pt,
    after skip=4pt,
    fontupper=\small
}

\newtcolorbox{definitionbox}{
    enhanced,
    breakable,
    colback=definitioncolor!5,
    colframe=definitioncolor,
    boxrule=0pt,
    leftrule=4pt,
    arc=0mm,
    left=6pt,
    right=6pt,
    top=4pt,
    bottom=4pt,
    before skip=4pt,
    after skip=4pt,
    fontupper=\small
}

\newtcolorbox{theorembox}{
    enhanced,
    breakable,
    colback=theoremcolor!5,
    colframe=theoremcolor,
    boxrule=0pt,
    leftrule=4pt,
    arc=0mm,
    left=6pt,
    right=6pt,
    top=4pt,
    bottom=4pt,
    before skip=4pt,
    after skip=4pt,
    fontupper=\small
}

\newtcolorbox{proofbox}{
    enhanced,
    breakable,
    colback=proofcolor!5,
    colframe=proofcolor,
    boxrule=0pt,
    leftrule=4pt,
    arc=0mm,
    left=6pt,
    right=6pt,
    top=4pt,
    bottom=4pt,
    before skip=4pt,
    after skip=4pt,
    fontupper=\small
}

% ---------- Команды ----------
\newcommand{\examquestion}[1]{
\begin{questionbox}
\textbf{Вопрос.} #1
\end{questionbox}
}

\newcommand{\examanswer}[1]{
\begin{answerbox}
\textbf{Ответ.} #1
\end{answerbox}
}

\newcommand{\examdefinition}[1]{
\begin{definitionbox}
\textbf{Определение.} #1
\end{definitionbox}
}

\newcommand{\examtheorem}[1]{
\begin{theorembox}
\textbf{Теорема.} #1
\end{theorembox}
}

\newcommand{\examproof}[1]{
\begin{proofbox}
\textbf{Доказательство.} #1
\end{proofbox}
}

\newcommand{\term}[1]{\textbf{#1}}
\newcommand{\important}[1]{\textit{#1}}

\newcommand{\examstep}[1]{
\par\smallskip
\noindent\textcolor{answercolor}{\rule{\linewidth}{0.4pt}}
\par\vspace{-2pt}
\noindent{\textbf{\textcolor{answercolor}{#1}}}
\par\vspace{-4pt}
\noindent\textcolor{answercolor}{\rule{\linewidth}{0.4pt}}
\par\smallskip
}

\begin{document}

\tableofcontents
\newpage

% =======================

\section{Дифференциальные уравнения}

\subsection{Линейные дифференциальные уравнения первого порядка}

\examquestion{Линейные дифференциальные уравнения первого порядка. Интегрирование линейных неоднородных дифференциальных уравнений первого порядка методом Бернулли (метод ``$u \cdot v$'') и методом Лагранжа (вариации произвольной постоянной). [Л. 15.]}

\examdefinition{
\term{Обыкновенным дифференциальным уравнением (ОДУ) первого порядка} называется уравнение, которое зависит от одной переменной $x$, неизвестной функции $y(x)$ и её производной:
\[
F(x,y(x),y'(x))=0.
\]
}

\examanswer{
Решение ЛНДУ методом Лагранжа (метод вариации произвольной постоянной)

Рассмотрим линейное неоднородное дифференциальное уравнение первого порядка:
\[
y'+p(x)y=f(x),
\]
где $p(x),f(x)$ — непрерывные функции на промежутке $I\subset \mathbb{R}$.

\examstep{Шаг 1. Решение соответствующего ЛОДУ}

Рассмотрим однородное уравнение:
\[
y'+p(x)y=0.
\]

Разделим переменные:
\[
\frac{dy}{dx}=-p(x)y.
\]

При $y\neq 0$:
\[
\frac{dy}{y}=-p(x)dx.
\]

Интегрируем:
\[
\int \frac{dy}{y}=-\int p(x)dx.
\]

Получаем:
\[
\ln|y|=-\int p(x)dx+C.
\]

Отсюда:
\[
y=Ce^{-\int p(x)dx}.
\]

Особое решение

При $y=0$:
\[
0+p(x)\cdot 0=0,
\]
то есть это тоже решение.

Оно входит в общее решение при $C=0$, поэтому отдельно его учитывать не нужно.

Общее решение ЛОДУ:
\[
y=Ce^{-\int p(x)dx}.
\]

\examstep{Шаг 2. Вариация постоянной}

Пусть теперь $C=C(x)$ — функция.

Ищем решение ЛНДУ в виде:
\[
y=C(x)e^{-\int p(x)dx}.
\]

\examstep{Шаг 3. Подстановка в исходное уравнение}

Найдём производную:
\[
y'=C'(x)e^{-\int p(x)dx}-C(x)p(x)e^{-\int p(x)dx}.
\]

Подставляем в уравнение:
\[
y'+p(x)y=f(x).
\]

Получаем:
\[
C'(x)e^{-\int p(x)dx}=f(x).
\]

\examstep{Шаг 4. Нахождение $C(x)$}

\[
C'(x)=f(x)e^{\int p(x)dx}.
\]

Интегрируем:
\[
C(x)=\int f(x)e^{\int p(x)dx}dx+k.
\]

\examstep{Итоговое решение ЛНДУ}

Подставляем $C(x)$ обратно:
\[
y=e^{-\int p(x)dx}\left(\int f(x)e^{\int p(x)dx}dx+k\right),\quad k=\mathrm{const}.
\]
}

\examanswer{
Решение ЛНДУ методом Бернулли (подстановка $y=u(x)v(x)$)

Рассмотрим линейное неоднородное уравнение:
\[
y'+p(x)y=f(x),
\]
где $p(x),f(x)$ — непрерывные функции на промежутке $I\subset \mathbb{R}$.

\examstep{Шаг 1. Подстановка}

Положим:
\[
y=u(x)v(x).
\]

Тогда:
\[
y'=u'v+uv'.
\]

Подставим в исходное уравнение:
\[
u'v+uv'+p(x)uv=f(x).
\]

Сгруппируем:
\[ 
v(u'+p(x)u) + (uv' -f(x)) = 0
\]

\examstep{Шаг 2. Выбор функции $u(x)$}

У нас есть свобода выбора. Удобно положить:
\[
u'+p(x)u=0.
\]

Это ЛОДУ.

\examstep{Шаг 3. Решение для $u(x)$}

\[
u'=-p(x)u. \text{- ДУ с раз. пер}
\]

Разделим переменные:
\[
\frac{du}{u}=-p(x)dx.
\]

Интегрируем:
\[
\ln|u|=-\int p(x)dx+C \quad \forall C = \mathrm{const}.
\]

\[
e^{\ln|u|} = e^{-\int p(x)\,dx + C}
\]

\[
u=Ce^{-\int p(x)dx}.
\]

Так как у нас имеется произвол по одной из переменных, то мы не только можем
выбрать ее, как нам удобно, но и конкретизировать общее решение (т.е. перейти от
общего решения к частному)

Для этого выберем $C = 1$ в последнем равенстве
\[
u=e^{-\int p(x)dx}.
\]

\examstep{Шаг 4. Подстановка в уравнение}

Подставляем в исходное выражение:
\[
u'v+uv'+p(x)uv=f(x).
\]

С учётом выбора $u$:
\[
uv'=f(x).
\]

\examstep{Шаг 5. Нахождение $v(x)$}

\[
v'=\frac{f(x)}{u}.
\]

Подставляем $u$:
\[
v'=f(x)e^{\int p(x)dx}.
\]

Интегрируем:
\[
v=\int f(x)e^{\int p(x)dx}dx+k.
\]

\examstep{Итоговое решение}

\[
y=u\cdot v=e^{-\int p(x)dx}\left(\int f(x)e^{\int p(x)dx}dx+k\right),\quad k=\mathrm{const}.
\]
}

\section*{Вопрос 22. Теорема Коши о существовании и единственности решения дифференциального уравнения $n$-го порядка. Интегрирование дифференциальных уравнений $n$-го порядка, допускающих понижение порядка}

\begin{theorembox}
\textbf{Теорема существования и единственности решения задачи Коши для дифференциального уравнения $n$-го порядка}

Пусть дано дифференциальное уравнение $n$-го порядка:
\[
y^{(n)} = f\left(x, y, y', \ldots, y^{(n-1)}\right).
\]

Если функция
\[
f\left(x, y, y', \ldots, y^{(n-1)}\right)
\]
и её частные производные
\[
\frac{\partial f}{\partial y},
\frac{\partial f}{\partial y'},
\ldots,
\frac{\partial f}{\partial y^{(n-1)}}
\]
непрерывны в некоторой области
\[
D \subset \mathbb{R}^{n+1},
\]
содержащей точку
\[
\left(x_0, y_0, y'_0, \ldots, y^{(n-1)}_0\right),
\]
то существует и притом единственное решение дифференциального уравнения $n$-го порядка, удовлетворяющее начальным условиям:
\[
y(x_0) = y_0,\quad
y'(x_0) = y'_0,\quad
\ldots,\quad
y^{(n-1)}(x_0) = y^{(n-1)}_0.
\]
\end{theorembox}

\begin{answerbox}
\textbf{Интегрирование дифференциальных уравнений $n$-го порядка, допускающих понижение порядка}

\textbf{1. Уравнения вида}
\[
y^{(n)} = f(x).
\]

Метод решения: $n$-кратное интегрирование:
\[
y = \int \ldots \int f(x)\,dx^n.
\]

При этом после каждого интегрирования появляется произвольная постоянная, поэтому общее решение содержит $n$ произвольных постоянных.

\[
y =
\underbrace{\int \ldots \int}_{n}
f(x)\,dx^n
+ C_1x^{n-1} + C_2x^{n-2} + \ldots + C_{n-1}x + C_n.
\]

\textbf{2. Уравнения, не содержащие переменную $x$ явно}

Рассмотрим уравнение вида:
\[
F\left(y, y', \ldots, y^{(n)}\right) = 0.
\]

Метод решения: вводится замена
\[
y' = p(y).
\]

Тогда производные выражаются через $p$ как функцию от $y$:
\[
y'' = \frac{dp}{dy}\cdot y' = p'(y)p,
\]
и далее:
\[
y^{(n)} = \text{выражение через } p,\ p',\ \ldots,\ p^{(n-1)}.
\]

После такой замены порядок уравнения понижается на единицу.

\textbf{3. Уравнения, не содержащие функцию $y$ в явном виде}

Рассмотрим уравнение вида:
\[
F\left(x, y', \ldots, y^{(n)}\right) = 0.
\]

Метод решения: вводится замена
\[
y' = p(x).
\]

Тогда
\[
y'' = p'(x),
\]
\[
y''' = p''(x),
\]
\[
\ldots
\]
\[
y^{(n)} = p^{(n-1)}(x).
\]

В результате исходное уравнение принимает вид:
\[
F\left(x, p, p', \ldots, p^{(n-1)}\right) = 0.
\]

То есть порядок уравнения понижается на единицу. После нахождения функции $p(x)$ решение $y(x)$ находится интегрированием:
\[
y = \int p(x)\,dx.
\]
\end{answerbox}

\end{document}